% !TEX root = main.tex
\section{Distribution of Multi-step Collatz Primitive Source}~\label{sec:distribution-of-multi-step-collatz-primitive-source}

Now let's consider the Collatz primitive source. They are those number that can't trace back again.
Since if $b \cdot \beta^m = \alpha\cdot a + \gamma$, then $m \in \Log_{\beta,\alpha}(b^{-1}\cdot\gamma)$.
But not all $b$ is reversible modulo $\alpha$, and not all $b^{-1} \cdot \gamma$ is a remainder of a power of $\beta$.

\begin{definition}[Collatz primitive source]~\label{collatz-primitive-source}
    For $b\in \widehat{\mathbb{N}}_{\beta}$, if there is a sequence $\bm{v}\in \mathbb{N}^{n}$ such that $f_{n}(\bm{v})$ is not reversible modulo $\alpha$, or
    \[\Log_{\beta, \alpha}(f_{n}(\bm{v})^{-1} \cdot \gamma) = \varnothing, \]
    then we say $\bm{v}$ is a $n$-step Collatz primitive source vector of $b$ and $f_{n}(\bm{v})$ is a $n$-step Collatz primitive source of $b$.
    For those $n = 0$, we simplly call them Collatz primitive source.
    Notice that a number $a$ is not reversible modulo $\alpha$ equivalents $\gcd(a, \alpha) > 1$. For these, we say they are trivial primitive source.
    Others we call them non-trivial primitive source.\newline
    And all non-trivial primitive source are denoted as 
    \[\tilde{\mathbb{S}} = \{b \in \widehat{\mathbb{N}}_{\beta} \mid \gcd(b, \alpha) = 1 \text{ and } \Log(b^{-1}\cdot\gamma) = \varnothing\}, \]
    all primitive source denotes as
    \[ \mathbb{S} = \{b \in \widehat{\mathbb{N}}_{\beta} \mid \gcd(b, \alpha) > 1\}\cup \tilde{\mathbb{S}}. \]
\end{definition}
\begin{proposition}~\label{prop:trivial-primitive-source}
    $\mathbb{S} - \tilde{\mathbb{S}} = \bigcup_{p\in\mathbb{P}}^{p\mid\alpha} p\widehat{\mathbb{N}}_{{\beta}}$.
\end{proposition}
\begin{proposition}~\label{prop:primitive-source-of-prime-alpha}
    If $\alpha \in \mathbb{P}$, then $\tilde{\mathbb{S}} = \varnothing$.
\end{proposition}
\begin{proposition}
    If $b \in \mathbb{S}$, then $\forall m \in \mathbb{Z}$, if $\beta \nmid (b + m\cdot\alpha)$, there is $b + m\cdot\alpha \in \mathbb{S}$.
\end{proposition}


But if $\beta \mid (b + m\cdot\alpha)$, then we have to turn it to $b' \equiv \frac{b + m\cdot\alpha}{\beta^{\Delta}}\pmod{\alpha}$.
This equivalence provides a new equation $b'\cdot \beta^{\Delta} = b + m'\cdot\alpha$, which is also a Collatz reduce formula, which $\gamma = b$.

From \cref{cor:modulo-alpha-of-layers}, we have
\begin{proposition}~\label{prop:all-source-value-set-has-primitive-source}
    Suppose $\gamma = 1$, then for $b \not\in\mathbb{S}$, we have $\mathbb{F}_{n}(b) \not\subset \mathbb{S}$ and $\mathbb{F}_{n}(b) \cap \mathbb{S} \ne \varnothing$.
\end{proposition}
% \begin{proof}
%     TODO.
% \end{proof}

For $(\alpha, \beta, \gamma) = (3,2,1)$, from \cref{prop:single-vary-rule} we have
\begin{proposition}
    Suppose $f_{n}(\bm{v}) \in \mathbb{S}$, then $\left\{f_{n}(\bm{v} + 6k\bm{e}_1) \mid k\in\mathbb{N}\right\} \subset \mathbb{S}$.
\end{proposition}